Students in higher education plan their studies against curricula that encode
prerequisites, electives and institutional rules, and the tools available to
them are split. Graph-based systems draw the structure of a curriculum but
leave the timeline aside; table-based planners arrange courses across semesters
and check the result against the regulations, but leave the structure implicit.
Where both have been brought into one tool, dependency information has been
embedded in the schedule itself, which ties it to a chronological layout.

This thesis designs, builds and evaluates a hybrid graph--table interface in
which a browsable curriculum graph and a schedule-managing table write to one
shared plan. The work proceeds in three steps. A formative study with six
students of two TU~Wien Informatics programmes, conducted around an interactive
but non-functional prototype, yields 64 themes that are translated into a ranked
specification of 13 features carrying 55 numbered requirements. From that
specification the tool is designed and implemented: a Graph View that draws the
curriculum as its containment hierarchy with formal prerequisites as a
switchable overlay, a Table View holding the same plan as semester lanes, a
Compliance Engine that checks the plan against two degree regulations after
every change, and a Recommendation Panel whose suggestions carry an explanation.
The tool is then evaluated with eleven students, each planning once without the
Graph View and once with it, observed through a frame-sampled record of the
interface rather than through self-report alone.

Planning took the shape of a single checklist-driven loop, the same for the most
fluent and the most junior participants. The graph did not extend that loop: it
took the place of the course catalogue at the step where candidates are found,
while courses continued to be committed in the table, and the period of the loop
was unchanged. Measuring the plans against their brief rather than against what
students reported revealed a calibration gap: every participant reported success
in runs where no plan met its brief and the tool reported nothing missing.

The evidence is descriptive. Two small samples at one institution, a fixed
scenario order, a single coder and descriptive statistics support an account of
how these students planned, not a claim that a component caused a behaviour.
What the work offers is a requirements specification traceable to student
evidence, a working artefact whose design decisions are recorded with their
costs, and a process-level account of planning with such a tool.
